\subsection{数据集描述}
\label{sec:dataset}

训练和评估依赖于在受控条件下获取的成对激光超声数据集。该数据集遵循成对采集协议：输入信号以有限次平均（N 次每位置）记录，参考目标采用高次平均（M 次每位置，其中 M $\gg$ N）以近似真值干净信号。本研究中使用的 N = 16 和 M = 256 的代表性值，使得输入与目标之间的标称信噪比(SNR)差异约为 12 dB；然而，确切的 N 和 M 值由实验硬件和采集时间约束决定。

\textbf{训练和验证试样。} 所有成对训练和验证数据均采集自一块完整铝板（300 mm $\times$ 300 mm $\times$ 3 mm）。选择铝作为试样材料是因为它是一种常见的航空航天结构合金，在超声频率下具有良好表征的兰姆波传播特性，能够可靠地解释声学特征（如到达时间和波形形态）。完整铝板提供了一个受控环境，用于建立具有一致噪声特性的成对数据，涵盖多次采集会话。

\textbf{缺陷测试试样。} 泛化性能在含机加工缺陷的铝试样上进行评估，包括贯穿厚度狭槽（宽度：1.0 mm；长度：20 mm）和边缘缺口（深度：1.5 mm；宽度：0.5 mm）。这些机加工缺陷在兰姆波响应中产生可测量的反射和模式转换，使其适用于评估去噪是否保留了诊断相关特征。独立的含缺陷试样仅用于测试——训练和验证阶段不使用任何缺陷数据。

\textbf{数据划分。} 完整铝板数据集按采集会话以 80/10/10 划分，得到约 80\% 的成对信号用于训练，10\% 用于验证（用于早停和超参数调优），10\% 用于保留测试以评估完整板性能。缺陷测试集由完全独立的试样组成，不属于划分比例的一部分。这确保了测试集提供从完整训练数据到含缺陷场景的跨域泛化的无偏评估。

\textbf{信号预处理。} 所有采集的信号均归一化至单位方差，并加窗以在 50 $\mu$s 门内隔离主兰姆波到达。具有可观察采集伪影（如激光功率波动或探头未对准）的信号在数据集组装前被丢弃，得到约 10,000 个成对信号实例的干净训练语料库，涵盖所有会话。